\subsection{Stochastic Knapsack Problem --- Paper Notes}\label{sec:sk-notes}

We summarize two papers on the stochastic knapsack problem that are relevant to the study of adaptivity gaps in stochastic combinatorial optimization.

Did not finish reading, especially the technical details. But would like to summarize the results first.

\subsubsection{DGV (2008)}

Consider a stochastic 0/1 knapsack with $n$ items. Each item $i$ has a deterministic value $v_i \ge 0$ and a random size $s_i$ drawn independently from a known distribution. The knapsack has capacity~1. Items are inserted sequentially; the process terminates when the knapsack overflows, and the overflowing item contributes no value. The goal is to maximize the expected total value of successfully placed items.

Two types of policies are considered:
\begin{itemize}
    \item \textbf{Nonadaptive:} commits to a fixed item ordering in advance.
    \item \textbf{Adaptive:} chooses the next item dynamically based on remaining capacity.
\end{itemize}

Here adaptivity immediately sounds helpful because we can get information about the current remaining capacity. This is impossible for nonadaptive algo. 

So what does adaptivity give us in stochastic TSP setting? We do not need the historical data, but we benefit from being able to change our mind once we probed some active vertices, which is impossible in nonadaptive algo. 

The \emph{adaptivity gap} is the worst-case ratio $\mathrm{ADAPT}(I) / \mathrm{NONADAPT}(I)$ over all instances $I$, measuring how much benefit comes from making decisions dynamically versus committing to a fixed ordering upfront. This is analogous to the integrality gap in integer programming.

\paragraph{Technical Tools}

\begin{itemize}
    \item \textbf{Effective value:} $w_i = v_i \cdot \P[s_i \le 1]$, an upper bound on the value any policy can obtain from item $i$.
    \item \textbf{Mean truncated size:} $\mu_i = \E[\min\{s_i, 1\}]$, which is not sensitive to the distribution in the range $s_i > 1$.
    \item \textbf{Martingale argument:} For any adaptive policy, the expected total truncated mass of items attempted satisfies $\E[\mu(A)] \le 2$ (Lemma~3.1 in~\cite{DeanGoemansVondrak2008}). This is the key structural lemma.
    \item \textbf{LP relaxation (simple):}
    \[
        \Phi(t) = \max\Big\{\sum_i w_i x_i \;:\; \sum_i \mu_i x_i \le t,\; x_i \in [0,1]\Big\}.
    \]
    Then $\mathrm{ADAPT} \le \Phi(2)$ (Theorem~3.1).
    \item \textbf{LP relaxation (stronger, polymatroid-based):} $\Psi(t)$ with exponential constraints capturing the mass bound more tightly. One has $\mathrm{ADAPT} \le \Psi(2) \le \Phi(2)$, and $\Psi(t)$ can be solved efficiently via a greedy algorithm on a polymatroid.
\end{itemize}

\paragraph{Main Results}

\begin{enumerate}
    \item \textbf{$\f{32}{7}$-approximation nonadaptive algorithm ($\approx 4.57$):}
    A randomized greedy algorithm that orders items by decreasing $w_i/\mu_i$, randomly selects one item to insert first, then follows the greedy order. Achieves expected value $\ge \f{7}{32} \cdot \mathrm{ADAPT}$. Can be derandomized.

    \item \textbf{$4$-approximation nonadaptive algorithm:}
    Uses the stronger polymatroid LP bound $\Psi(2)$. A simple greedy nonadaptive policy achieves expected value $\ge \f{1}{4} \cdot \mathrm{ADAPT}$.

    \item \textbf{$(3+\eps)$-approximation adaptive algorithm:}
    A polynomial-time adaptive policy combining a $(2+\eps)$-approximation nonadaptive policy for ``small'' items with a $(1+\eps)$-approximation adaptive policy for ``large'' items.

    \item \textbf{Adaptivity gap $\le 4$} for Non-Risky stochastic knapsack.

    \item \textbf{Ordered adaptive model:} Items must be processed in a given order; the decision to insert or skip is adaptive.
    \begin{itemize}
        \item Optimal ordered adaptive $\le 4 \cdot \mathrm{NONADAPT}$ for any ordering.
        \item Optimal ordered adaptive $\le 9.5 \cdot$ optimal fully adaptive for the worst-case ordering.
    \end{itemize}
\end{enumerate}

\paragraph{Key Examples}

\begin{itemize}
    \item \textbf{Example~1:} A knapsack with $n$ unit-value items demonstrates that an optimal adaptive policy can have exponentially many leaves in its decision tree.
    \item \textbf{Example~3:} Shows the simple LP bound $\Phi(2)$ can be off by a factor of $\sim 4$ from $\mathrm{ADAPT}$, so $\f{32}{7} \approx 4.57$ cannot be improved with Theorem~3.1 alone.
    \item \textbf{Example~5:} The randomized greedy algorithm can be as bad as $\mathrm{ADAPT}/4$ on instances with items of size $(1+\eps)/2$ versus Bernoulli items.
\end{itemize}

\subsubsection{Barak and Talgam-Cohen (2026): Semi-Adaptivity Gaps and Improved Approximation}

Available as arXiv:2602.24042v2, March 2026~\cite{BarakTalgamCohen2026}.

\paragraph{Problem Setup}

The same stochastic knapsack framework as Dean et al., but with two variants:
\begin{itemize}
    \item \textbf{Non-Risky-SK:} Overflow loses only the overflowing item's value.
    \item \textbf{Risky-SK (main focus):} Overflow forfeits \emph{all} accumulated value. Also called Fixed-Set, All-Or-None, or Blackjack variant.
\end{itemize}

\paragraph{Key New Concept --- Semi-Adaptivity}

A \emph{$k$-semi-adaptive} policy makes at most $k$ adaptive queries (observations of remaining knapsack capacity). Between queries, items are inserted nonadaptively.
\begin{itemize}
    \item $k = 0$: fully nonadaptive.
    \item $k = n$: fully adaptive.
\end{itemize}
This provides a smooth interpolation between the two extremes.

\paragraph{Two New Gap Notions}

\begin{itemize}
    \item \textbf{$0$-$k$ semi-adaptivity gap} ($G_{0\text{-}k}$): the ratio $A_k(I) / \mathrm{ALG}(I)$, measuring the benefit of augmenting a nonadaptive policy with $k$ queries.
    \item \textbf{$k$-$n$ semi-adaptivity gap} ($G_{k\text{-}n}$): the ratio $\mathrm{ADAPT}(I) / A_k(I)$, measuring the cost of limiting to $k$ queries instead of full adaptivity.
    \item \textbf{Decomposition (Lemma~1):} $G_{k\text{-}n} \ge G / G_{0\text{-}k}$, i.e., the full adaptivity gap factors through the semi-adaptivity gaps.
\end{itemize}

\paragraph{Main Results --- Risky-SK}

\begin{center}
\begin{tabular}{l c c c c}
    \hline
    Quantity & Bound & Their Result & Previous Best & Location \\
    \hline
    $0$-$n$ full adaptivity gap & Upper & $8.47$ & $9.5$~\cite{DeanGoemansVondrak2008} & Thm.~1, Sec.~4 \\
    $0$-$n$ full adaptivity gap & Lower & $2$ & $1.5$~\cite{LV14} & Thm.~8, Sec.~7 \\
    $0$-$1$ semi-adaptivity gap & Upper & $1.69$ & --- & Thm.~7, Sec.~6 \\
    $0$-$1$ semi-adaptivity gap & Lower & $1.69$ & --- & Thm.~7, Sec.~6 \\
    $1$-$n$ semi-adaptivity gap & Upper & $8.26$ & --- & Thm.~2, Sec.~4 \\
    $1$-$n$ semi-adaptivity gap & Lower & $1.18$ & --- & Cor.~3, Sec.~7 \\
    $k$-$n$ gap, $k = \tilde{O}(1/\eps)$ & Upper & $6.44 + \sqrt{\eps}$ & --- & Thm.~4, Sec.~5 \\
    \hline
\end{tabular}
\end{center}

\paragraph{Main Results --- Non-Risky-SK}

\begin{center}
\begin{tabular}{l c c c c}
    \hline
    Quantity & Bound & Their Result & Previous Best & Location \\
    \hline
    $0$-$n$ gap ($\eps$-Noisy Bernoulli) & Upper & $2$ & $4$~\cite{DeanGoemansVondrak2008} & Prop.~5, Sec.~E \\
    $0$-$n$ full adaptivity gap & Lower & $1.37$ & $1.25$~\cite{DeanGoemansVondrak2008} & Cor.~4, Sec.~E \\
    $k$-$n$ gap, $k = \tilde{O}(1/\eps)$ & Upper & $3 + \eps$ & --- & Thm.~3, Sec.~5 \\
    \hline
\end{tabular}
\end{center}

\paragraph{Algorithmic Improvements over DGV08}

\begin{center}
\resizebox{\textwidth}{!}{%
\begin{tabular}{l c c}
    \hline
    & DGV08 & Barak \& Talgam-Cohen \\
    \hline
    Nonadaptive algorithm (Risky-SK) & $9.5$-approx & $8.47$-approx (Alg.~1) \\
    Running time & $O(\mathrm{poly}(n))$, $\eps$-dependent & $O(n \log n)$ \\
    Best known (Risky-SK) & $(8+\eps)$-approx fully adaptive & $6.44$-approx, $\tilde{O}(1/\eps)$ queries \\
    \hline
\end{tabular}%
}
\end{center}

\paragraph{Key Insights}

\begin{enumerate}
    \item \textbf{A single adaptive query matters.} The $0$-$1$ semi-adaptivity gap is exactly $1 + \ln 2 \approx 1.69$ for Risky-SK. Even one observation of remaining capacity gives a provable constant-factor improvement.

    \item \textbf{Constant queries $\approx$ full adaptivity.} With $\tilde{O}(1/\eps)$ adaptive queries, the $k$-$n$ gap is only $6.44 + \sqrt{\eps}$, improving on the best known fully adaptive $8$-approximation~\cite{FLX18}.

    \item \textbf{Risky-SK requires optimal stopping.} Unlike Non-Risky-SK where it is always optimal to keep inserting until overflow, Risky-SK policies must balance the gain of adding an item against the risk of losing everything. The optimal Risky-SK policy has non-negligible overflow probability ($\approx 0.37$).

    \item \textbf{$\eps$-Noisy Bernoulli distributions} (size $\eps$ with probability $1-\eps$, size $1$ with probability $\eps$) form a natural hard family: the adaptivity gap lower bound of $2$ for Risky-SK and the tight $0$-$1$ gap bound of $1.69$ are both established on this family.
\end{enumerate}

\paragraph{Technique --- Simplify-Equalize-Optimize}

A three-step approach for analyzing adaptive decision trees:
\begin{enumerate}
    \item \textbf{Simplify} the decision tree structure.
    \item \textbf{Equalize} parameters across branches.
    \item \textbf{Optimize} the simplified, equalized tree.
\end{enumerate}
This technique may have broader applications to semi-adaptivity in other stochastic combinatorial optimization problems.

\subsubsection{Relationship Between the Two Papers}

\begin{itemize}
    \item Barak and Talgam-Cohen directly build on and improve the Dean et al.\ framework.
    \item The core LP relaxation and martingale-based analysis from~\cite{DeanGoemansVondrak2008} remain foundational.
    \item The new semi-adaptivity framework provides a finer-grained understanding of when and how much adaptivity helps.
    \item All improved bounds tighten the gap between upper and lower bounds, but significant gaps remain (e.g., Risky-SK: lower bound $2$ vs.\ upper bound $8.47$).
\end{itemize}
